
\documentclass[12pt]{amsart}
\usepackage{geometry} % see geometry.pdf on how to lay out the page. There's lots.
\geometry{a4paper} % or letter or a5paper or ... etc
\usepackage[T1]{fontenc}
\usepackage[latin9]{inputenc}
\usepackage{amsmath}
\usepackage{amsaddr}
\usepackage{hyperref}
\usepackage{dirtytalk}
\usepackage{float}
\usepackage{listings}
\usepackage{color}
\usepackage{tikz}
\usepackage{enumitem}

 
\definecolor{codegreen}{rgb}{0,0.6,0}
\definecolor{codegray}{rgb}{0.5,0.5,0.5}
\definecolor{stringcolor}{rgb}{0.7,0.23,0.36}
\definecolor{backcolour}{rgb}{0.95,0.95,0.92}
\definecolor{keycolor}{rgb}{0.007,0.01,1.0}
\definecolor{itemcolor}{rgb}{0.01,0.0,0.49}
 
\lstdefinestyle{mystyle}{
    %backgroundcolor=\color{backcolour},   
    commentstyle=\color{codegreen},
    keywordstyle=\color{keycolor},
    numberstyle=\tiny\color{codegray},
    stringstyle=\color{stringcolor},
    basicstyle=\footnotesize,
    breakatwhitespace=false,         
    breaklines=true,                 
    captionpos=b,                    
    keepspaces=true,                 
    numbers=left,                    
    numbersep=5pt,                  
    showspaces=false,                
    showstringspaces=false,
    showtabs=false,                  
    tabsize=2
}
 
\lstset{style=mystyle}

\lstdefinelanguage{Swift}{
  keywords={associatedtype, class, deinit, enum, extension, func, import, init, inout, internal, let, operator, private, protocol, public, static, struct, subscript, typealias, var, break, case, continue, default, defer, do, else, fallthrough, for, guard, if, in, repeat, return, switch, where, while, as, catch, dynamicType, false, is, nil, rethrows, super, self, Self, throw, throws, true, try, associativity, convenience, dynamic, didSet, final, get, infix, indirect, lazy, left, mutating, none, nonmutating, optional, override, postfix, precedence, prefix, Protocol, required, right, set, Type, unowned, weak, willSet},
  ndkeywords={class, export, boolean, throw, implements, import, this},
  sensitive=false,
  comment=[l]{//},
  morecomment=[s]{/*}{*/},
  morestring=[b]',
  morestring=[b]"
}

\lstset{emph={Int,count,abs,repeating,Array}, emphstyle=\color{itemcolor}}


\title{Week 08}

\date{\today}

\lstset{style=mystyle}

%%% BEGIN DOCUMENT
\begin{document}
\maketitle


Name: 
\\

Collaborators (if any): 

\section{Tasks}

\subsection{True or False}
Choose T or F for each of the following statements and briefly explain why.  
\begin{enumerate}
\item T/F In the potential method for amortized analysis, the potential energy
should never go negative. 
\item T/F The quicksort algorithm that uses linear-time median finding to run in
worst-case $O(n \: log \: n)$ time requires $\Theta (n)$ auxiliary space.
\end{enumerate}

\section{Exercises/Problems}

\subsection{Amortized Analysis} 
Design a data structure to maintain a set S of n distinct integers that supports the following two operations:
\\

\begin{enumerate} 
\item $Insert(x, S)$ : insert integer $x$ into $S$
\item $Remove.Bottom.Half(S)$ : remove the smallest $\lceil\frac{n}{2}\rceil$ integers from $S$
\\
\end{enumerate}

Describe your algorithm and give the worse-case time complexity of the two operations. Then
carry out an amortized analysis to make Insert(x, S) run in amortized O(1) time, and Remove.Bottom.Half(S) run in amortized 0 time. 

\subsection{Quicksort} 
Write up some code for a quicksort algorithm. Be sure to specify which pivot selection method you are using (reminder on pivot selection choices from lecture: always pick the fist element as pivot, always pick the lat element as pivot, pick the median as pivot, or pick a random pivot). Code should be clearly commented to provide context and steps taken. 
\end{document}
